\documentclass[aspectratio=169]{beamer}

\usetheme{Madrid}
\useinnertheme{circles}
\usecolortheme{default}

\usepackage[utf8]{inputenc}
\usepackage[T5]{fontenc}
\usepackage[vietnamese]{babel}
\usepackage{amsmath}
\usepackage{amssymb}
\usepackage{booktabs}
\usepackage{array}
\usepackage{tikz}
\usetikzlibrary{arrows.meta, positioning}

\definecolor{donegreen}{RGB}{29,125,75}
\definecolor{partialorange}{RGB}{196,112,21}
\definecolor{notred}{RGB}{176,49,49}
\definecolor{softblue}{RGB}{34,86,166}

\hypersetup{
  pdftitle={Hệ thống kiểm soát ra vào thông minh - Báo cáo cuối kỳ},
  pdfauthor={Nhóm: Quang, Nhân, Vân, Phú},
  pdfsubject={Đánh giá sản phẩm AIoT với Camera, RFID và Web Sentinel},
}

\title[Kiểm soát ra vào AIoT]{Hệ thống kiểm soát ra vào thông minh}
\subtitle{Báo cáo cuối kỳ \& Đánh giá mức độ đáp ứng mục tiêu}
\author[Nhóm thực hiện]{Lưu Huy Minh Quang -- 23127016 \and Hoàng Nhân -- 23127094 \and Trần Cao Vân -- 23127141 \and Nguyễn Trần Thiên Phú -- 23127246}
\institute[HCMUS]{Khoa Công nghệ Thông tin\\Đại học Khoa học Tự nhiên, ĐHQG TP.HCM\\[0.3em]\textit{Môn: Nhập môn lập trình điều khiển thiết bị thông minh}}
\date{Tháng 8, 2026}

\setbeamertemplate{navigation symbols}{}

\begin{document}

%------- TITLE -------
\begin{frame}
  \begin{center}
    \vspace{0.3em}
    {\usebeamerfont{title}\Large\bfseries\inserttitle\par}
    \vspace{0.3em}
    {\usebeamerfont{subtitle}\insertsubtitle\par}

    \vspace{1.0em}
    {\large\bfseries Nhóm thực hiện}\\[0.5em]
    \begin{tabular}{ll}
      Lưu Huy Minh Quang & 23127016 \\
      Hoàng Nhân & 23127094 \\
      Trần Cao Vân & 23127141 \\
      Nguyễn Trần Thiên Phú & 23127246 \\
    \end{tabular}\\[0.7em]
    \textsl{Khoa CNTT -- ĐH Khoa học Tự nhiên, ĐHQG TP.HCM}\\
    \textsl{Môn: Nhập môn lập trình điều khiển thiết bị thông minh}

    \vspace{0.7em}
    {\large \insertdate}
  \end{center}
\end{frame}

%------- NỘI DUNG -------
\begin{frame}{Nội dung trình bày}
  \vspace{0.3em}
  \begin{enumerate}
    \item Sản phẩm cuối cùng đang làm được gì?
    \item Các phần từng ghi chưa hoàn thành nay ở trạng thái nào?
    \item Đồ án đáp ứng ba mục tiêu ban đầu đến đâu?
    \item Kết luận và phạm vi của phiên bản hiện tại
  \end{enumerate}
\end{frame}

\section{Kết quả sản phẩm}

\begin{frame}{Kết luận sau khi rà soát toàn bộ đồ án}
  \begin{block}{Sản phẩm đã thành hình ở mức nguyên mẫu một cổng}
    Người dùng có thể đăng ký hồ sơ, dùng khuôn mặt hoặc thẻ để xác thực, mở cổng tự động và theo dõi sự kiện trên trang quản lý.
  \end{block}

  \vspace{0.3em}
  \begin{itemize}
    \item \textcolor{donegreen}{\textbf{Đã có:}} luồng nhận diện khuôn mặt, xác thực thẻ, điều khiển cổng, lưu hồ sơ và gửi cảnh báo.
    \item \textcolor{donegreen}{\textbf{Đã phát hiện được hành vi trèo:}} cảm biến siêu âm báo vi phạm khi có vật thể trong vùng kiểm soát lúc cổng chưa mở.
    \item \textcolor{partialorange}{\textbf{Nhận diện còn chậm:}} hệ thống nhận diện đúng luồng, nhưng toàn bộ quá trình mất khoảng 5--7 giây.
  \end{itemize}
\end{frame}

\begin{frame}{Sản phẩm cuối cùng đang làm được gì?}
  \begin{columns}[T]
    \column{0.48\textwidth}
      \textbf{Tại cổng}
      \begin{itemize}
        \item Phát hiện có người tiến vào vùng cổng và phát hiện trèo khi cổng chưa mở.
        \item Nhận diện khuôn mặt đã đăng ký ngay trên camera.
        \item Đọc thẻ, phân biệt thẻ hợp lệ và không hợp lệ.
        \item Mở/đóng thanh chắn, đổi đèn báo và phát tiếng xác nhận.
      \end{itemize}

    \column{0.48\textwidth}
      \textbf{Trên trang quản lý}
      \begin{itemize}
        \item Quản lý hồ sơ, vai trò, email và thẻ của người dùng.
        \item Đăng ký khuôn mặt theo ba góc nhìn.
      \item Xem trực tiếp camera, lịch sử xác thực và cảnh báo gần nhất.
        \item Gửi email đúng nhóm quản trị hoặc an ninh khi có cảnh báo.
      \end{itemize}
  \end{columns}
\end{frame}

\begin{frame}{Một lượt ra vào diễn ra như thế nào?}
  \begin{center}
  \begin{tikzpicture}[
    flowstep/.style={rectangle, rounded corners=4pt, draw=softblue, fill=blue!7,
      text width=2.25cm, align=center, minimum height=1.15cm, font=\small\bfseries},
    arr/.style={-{Stealth[length=2mm]}, thick, softblue}
  ]
    \node[flowstep] (sense) {Có người đến cổng};
    \node[flowstep, right=0.55cm of sense] (auth) {Nhận diện mặt hoặc đọc thẻ};
    \node[flowstep, right=0.55cm of auth] (decide) {Kiểm tra quyền ra vào};
    \node[flowstep, right=0.55cm of decide] (act) {Mở cổng hoặc giữ khóa};
    \node[flowstep, right=0.55cm of act] (record) {Ghi lịch sử và cảnh báo};

    \draw[arr] (sense) -- (auth);
    \draw[arr] (auth) -- (decide);
    \draw[arr] (decide) -- (act);
    \draw[arr] (act) -- (record);
  \end{tikzpicture}
  \end{center}

  \vspace{0.7em}
  \begin{block}{Điểm thay đổi so với thiết kế cũ}
    Dashboard có thể xem trực tiếp luồng camera tại cổng. Việc nhận diện vẫn được xử lý ngay trên camera và kết quả xác thực được gửi về mạch cổng.
  \end{block}
\end{frame}

\section{Đối chiếu phần chưa hoàn thành}

\begin{frame}{Các phần từng ghi chưa hoàn thành: kết quả cuối (1/2)}
  \begin{itemize}
    \item \textcolor{donegreen}{\textbf{Đăng ký và nhận diện khuôn mặt -- Đã có chức năng:}}
      camera nhận ba góc mặt, lưu dữ liệu ngay trên thiết bị, nhận diện và gửi kết quả về cổng.
    \item \textcolor{donegreen}{\textbf{Phân nhóm email -- Đã có chức năng:}}
      cảnh báo được gửi tới người có vai trò Quản trị viên hoặc Nhân viên an ninh; email cá nhân dùng cho thông báo lượt vào hợp lệ.
    \item \textcolor{donegreen}{\textbf{Luồng camera -- Đã có chức năng:}}
      dashboard có thể xem trực tiếp luồng hình ảnh từ camera; việc nhận diện vẫn diễn ra tại camera và kết quả được gửi về cổng.
  \end{itemize}
\end{frame}

\begin{frame}{Các phần từng ghi chưa hoàn thành: kết quả cuối (2/2)}
  \begin{itemize}
     \item \textcolor{donegreen}{\textbf{Phát hiện trèo -- Đã có chức năng:}}
      khi cổng chưa mở mà cảm biến siêu âm phát hiện vật thể trong vùng kiểm soát, hệ thống xem đây là vi phạm, tiếp tục khóa cổng và gửi cảnh báo.
    \item \textcolor{donegreen}{\textbf{Khôi phục kết nối -- Đã có cơ chế tự thử lại:}}
      mạch và trang quản lý tự nối lại khi Wi-Fi hoặc kênh truyền bị gián đoạn ngắn.
  \end{itemize}
\end{frame}

\section{Đánh giá ba mục tiêu ban đầu}

\begin{frame}{Mục tiêu 1: Đã nhận diện được, nhưng phản hồi còn chậm}
  \begin{block}{Mục tiêu ban đầu}
    Nhận diện tối đa 20 người trong phạm vi 0,5 m; độ chính xác trên 95\%, nhận nhầm dưới 2\% và phản hồi dưới 1,5 giây.
  \end{block}

  \begin{itemize}
    \item \textbf{Đã đáp ứng về chức năng:} đăng ký ba góc mặt, lưu dữ liệu trên camera, tự nhận diện và mở cổng khi khớp.
    \item \textbf{Phù hợp quy mô nguyên mẫu:} hệ thống quản lý được hồ sơ người dùng và giới hạn tối đa 20 thẻ được đồng bộ xuống cổng.
    \item \textbf{Thời gian thực tế:} từ lúc phát hiện người đến khi nhận diện và phản hồi mất khoảng 5--7 giây.
  \end{itemize}

  \begin{alertblock}{Đánh giá: Đạt chức năng, chưa đạt mục tiêu tốc độ}
    Hệ thống đã phát hiện và nhận diện được người dùng, nhưng phản hồi chậm hơn mức 1,5 giây đặt ra ban đầu.
  \end{alertblock}
\end{frame}

\begin{frame}{Mục tiêu 2: Đã phát hiện trèo bằng cảm biến siêu âm}
  \begin{block}{Mục tiêu ban đầu}
    Phát hiện hành vi trèo hoặc vượt cổng khi chưa được cấp quyền, giữ cổng ở trạng thái an toàn và gửi cảnh báo cho người quản lý.
  \end{block}

  \begin{itemize}
    \item \textbf{Điều kiện phát hiện:} cảm biến siêu âm ghi nhận vật thể trong vùng kiểm soát trong khi cổng vẫn đóng và chưa có lệnh mở hợp lệ.
    \item \textbf{Cách xử lý:} hệ thống xem đây là hành vi trèo/vượt cổng, tiếp tục giữ khóa, bật trạng thái cảnh báo và gửi sự kiện lên trang quản lý.
    \item \textbf{Tránh báo nhầm lượt hợp lệ:} khi người dùng đã xác thực và cổng được mở, tín hiệu đi qua không bị xem là hành vi trèo.
  \end{itemize}

  \begin{alertblock}{Đánh giá: Đạt mục tiêu phát hiện và cảnh báo}
    Cơ chế siêu âm kết hợp trạng thái đóng/mở của cổng đã thực hiện được việc phát hiện trèo theo thiết kế của nhóm.
  \end{alertblock}
\end{frame}

\begin{frame}{Mục tiêu 3: Luồng RFID/NFC đã hoàn chỉnh}
  \begin{block}{Mục tiêu ban đầu}
    Xác thực RFID/NFC và mở cổng dưới 2 giây, độ chính xác tối thiểu 98\% trong điều kiện bình thường.
  \end{block}

  \begin{itemize}
    \item \textbf{Đã đáp ứng về chức năng:} quét và gán thẻ, nhận thẻ hợp lệ/không hợp lệ, mở cổng, đổi đèn, phát tiếng báo và ghi lịch sử.
    \item \textbf{Có phương án dự phòng:} danh sách thẻ được lưu ngay trên mạch nên thẻ hợp lệ có thể được kiểm tra nhanh mà không phải chờ máy chủ trả lời.
    \item \textbf{Trải nghiệm sử dụng:} thẻ hợp lệ được xử lý ngay tại cổng; thẻ không hợp lệ bị từ chối và ghi lại trong lịch sử.
  \end{itemize}

  \begin{alertblock}{Đánh giá: Đạt mục tiêu chức năng}
    Luồng từ quét thẻ đến mở cổng tự động đã hoạt động đầy đủ trong sản phẩm.
  \end{alertblock}
\end{frame}

\begin{frame}{Tổng hợp mức độ đáp ứng ba mục tiêu}
  \begin{center}
  \renewcommand{\arraystretch}{1.45}
  \begin{tabular}{>{\raggedright\arraybackslash}p{2.2cm} >{\raggedright\arraybackslash}p{5.1cm} >{\raggedright\arraybackslash}p{3.3cm}}
    \toprule
    \textbf{Mục tiêu} & \textbf{Kết quả sản phẩm} & \textbf{Mức đáp ứng} \\
    \midrule
    Nhận diện người & Đã phát hiện, nhận diện và mở cổng; toàn bộ quá trình mất khoảng 5--7 giây & \textcolor{partialorange}{\textbf{Đạt chức năng, còn chậm}} \\
    Phát hiện trèo/nhảy & Cảm biến phát hiện vật thể khi cổng chưa mở, giữ khóa và gửi cảnh báo & \textcolor{donegreen}{\textbf{Đạt}} \\
    RFID/NFC & Luồng quét thẻ, kiểm tra quyền, mở cổng và ghi lịch sử đã hoạt động & \textcolor{donegreen}{\textbf{Đạt}} \\
    \bottomrule
  \end{tabular}
  \end{center}
\end{frame}

\section{Kiến trúc và kết luận}

\begin{frame}{Kiến trúc sản phẩm cuối cùng}
  \begin{center}
  \begin{tikzpicture}[
    layerbox/.style={rectangle, rounded corners=4pt, text width=2.85cm,
      align=center, minimum height=1.25cm, font=\small\bfseries, inner sep=5pt},
    inputbox/.style={layerbox, draw=blue!70, fill=blue!8},
    ctrlbox/.style={layerbox, draw=violet!70!black, fill=violet!10},
    actbox/.style={layerbox, draw=green!55!black, fill=green!9},
    webbox/.style={layerbox, draw=orange!80!black, fill=orange!10},
    arr/.style={-{Stealth[length=2mm]}, thick, blue!70},
    dataarr/.style={<->, >=Stealth, thick, orange!80!black}
  ]
    \node[inputbox] (input) {Thiết bị đầu vào\\{\scriptsize Camera + thẻ + cảm biến}};
    \node[ctrlbox, right=0.8cm of input] (ctrl) {Mạch điều khiển cổng\\{\scriptsize Kiểm tra quyền ra vào}};
    \node[actbox, right=0.8cm of ctrl] (act) {Thiết bị phản hồi\\{\scriptsize Servo + đèn + còi}};
    \node[webbox, below=1.0cm of ctrl] (web) {Trang quản lý\\{\scriptsize Hồ sơ + lịch sử + email}};

    \draw[arr] (input) -- (ctrl);
    \draw[arr] (ctrl) -- (act);
    \draw[dataarr] (ctrl) -- node[left, font=\scriptsize]{Trạng thái / lệnh đăng ký} (web);
  \end{tikzpicture}
  \end{center}

  \vspace{0.4em}
  \small\textbf{Nguyên tắc chính:} việc nhận diện được xử lý ngay tại camera; mạch cổng quyết định mở/đóng; trang web dùng để quản lý hồ sơ, theo dõi và nhận cảnh báo.
\end{frame}

\begin{frame}{Phạm vi của phiên bản cuối}
  \begin{columns}[T]
    \column{0.48\textwidth}
      \textbf{Có thể trình diễn}
      \begin{itemize}
        \item Đăng ký hồ sơ, khuôn mặt và thẻ.
        \item Mở cổng bằng khuôn mặt hoặc RFID.
        \item Phát hiện trèo khi cảm biến có tín hiệu lúc cổng chưa mở.
        \item Từ chối người/thẻ không hợp lệ.
        \item Hiển thị lịch sử và gửi cảnh báo.
      \end{itemize}

    \column{0.48\textwidth}
      \textbf{Giới hạn còn lại}
      \begin{itemize}
        \item Nhận diện khuôn mặt cần khoảng 5--7 giây cho toàn bộ quá trình.
        \item Nhật ký chủ yếu được giữ trên trình duyệt đang sử dụng.
      \end{itemize}
  \end{columns}
\end{frame}

\begin{frame}{Kết luận}
  \begin{block}{Kết quả cuối kỳ}
    Nhóm đã xây dựng được một nguyên mẫu kiểm soát ra vào có đủ luồng người dùng, thiết bị và trang quản lý để trình diễn. Điểm mạnh nhất là xác thực RFID và nhận diện khuôn mặt ngay trên thiết bị.
  \end{block}

  \vspace{0.4em}
  \begin{alertblock}{Kết luận so với ba mục tiêu}
    Mục tiêu 2 và 3 đã được đáp ứng trong sản phẩm. Mục tiêu 1 đã phát hiện và nhận diện được người dùng, nhưng thời gian toàn bộ quá trình còn khoảng 5--7 giây, chậm hơn mục tiêu ban đầu.
  \end{alertblock}

  \vspace{0.5em}
  \begin{center}
    {\Large\bfseries Q\&A -- Phản hồi và thảo luận}
  \end{center}
\end{frame}

\end{document}
